\chapter{Introduction}
\label{ch:introduction}

\section{Motivation}

Gliomas are the most common primary tumours of the central nervous system, and
glioblastoma (GBM, WHO grade~IV) is their most aggressive and most frequent
malignant form in adults. Despite advances in surgery, radiotherapy and
chemotherapy, the prognosis remains poor, with a median survival on the order of
only 12--15 months after diagnosis. Because glioblastoma is highly infiltrative
and heterogeneous, its early and precise characterisation is directly relevant to
treatment planning, surgical guidance, radiotherapy target definition and the
monitoring of disease progression.

Magnetic resonance imaging (MRI) is the primary modality for assessing brain
tumours because it provides excellent soft-tissue contrast without ionising
radiation. In clinical practice several complementary sequences are acquired: a
native T1-weighted scan (T1), a contrast-enhanced (gadolinium) T1-weighted scan
(T1Gd), a T2-weighted scan (T2) and a T2 fluid-attenuated inversion-recovery scan
(FLAIR). Each sequence highlights different tissue properties: the enhancing
tumour rim is most visible on T1Gd, while peritumoral edema and infiltration are
best appreciated on T2 and FLAIR. A complete assessment therefore requires the
radiologist to reason jointly across all four sequences.

\section{Problem statement}

Delineating a glioblastoma and its sub-regions by hand is time-consuming and
requires expert knowledge. A single case contains on the order of a hundred and
fifty axial slices per sequence, and manual annotation of the tumour core,
enhancing tumour and surrounding edema across the full volume can take a
radiologist a considerable amount of time. Manual segmentation also suffers from
substantial inter- and intra-rater variability, which limits reproducibility and
complicates quantitative follow-up.

Automating this task is difficult. Tumours vary widely in size, shape and
location; intensities are not standardised across scanners and acquisitions; and
the boundaries between sub-regions such as edema and healthy tissue are often
diffuse. An automated method must integrate information from all four modalities,
respect the three-dimensional structure of the tumour, and produce a dense
voxel-level labelling rather than a single image-level decision.

\section{Objectives and contributions}

The objective of this project is to build a reliable, reproducible and fully
automated pipeline that segments glioblastoma sub-regions from multi-parametric
MRI, and to evaluate it honestly on a public benchmark. The specific
contributions are:

\begin{itemize}
  \item A complete, reproducible training and evaluation pipeline for 3D brain
        tumour segmentation, built on PyTorch and MONAI and configured through a
        single entry point, with fixed seeds and a deterministic data split.
  \item A study focused specifically on \emph{glioblastoma}: the BraTS~2020
        cases are filtered by tumour grade so that the model is trained and
        evaluated on the 293 high-grade-glioma (glioblastoma) subjects, keeping
        the ``glioblastoma'' framing of this work accurate rather than mixing in
        lower-grade gliomas.
  \item A quantitative evaluation using both an overlap metric (Dice similarity
        coefficient) and a boundary metric (95th-percentile Hausdorff distance),
        reported per sub-region, together with qualitative overlays that place the
        prediction next to the ground truth.
  \item An honest discussion of the model's behaviour, its failure cases and its
        limitations, and a concrete path towards a deployable web application
        that detects and localises tumours from an uploaded scan.
\end{itemize}

\section{Report structure}

Chapter~\ref{ch:related} reviews prior work on brain tumour segmentation, from
classical methods to modern encoder--decoder and transformer architectures.
Chapter~\ref{ch:method} describes the dataset, the pre-processing, the 3D~U-Net
architecture, the loss function and the training and evaluation protocol.
Chapter~\ref{ch:results} reports the quantitative and qualitative results.
Chapter~\ref{ch:discussion} interprets these results and states the limitations
of the study, and Chapter~\ref{ch:conclusion} concludes and outlines future work.
